% ======================================================================
%  chapitre2.tex — Implementation
%  \input from main.tex — do not compile standalone.
% ======================================================================

\chapter{Implementation}
\label{chap:implementation}

This chapter describes the concrete implementation of the translation methodology presented in \autoref{chap:translation}. The translator (\texttt{archimate\_to\_eakg.py}) converts an ArchiMate Model Exchange File Format XML document into an RDF/OWL knowledge graph. Enrichment is performed directly in the triplestore (GraphDB) by loading the ontologies and inserting necessary triples using SPARQL queries. The entire stack is built on open standards and open-source tooling. A YARRRML implementation of the translation rules was provided as a contribution by Amel BELMAKSENE to the work.

% ------------------------------------------------------------------
\section{Architecture of the Translator}
\label{sec:translator}

\begin{figure}[H]
    \centering
    \includegraphics[width=\linewidth]{implementation_updated.png}
    \caption{Translator and enrichment pipeline}
    \label{fig:translator-enrichment}
\end{figure}

The translator (\texttt{archimate\_to\_eakg.py}) takes an ArchiMate Model Exchange File Format (v3.1) XML file as input and produces an Enterprise Architecture Knowledge Graph serialized as a Turtle\parencite{turtle} (\texttt{.ttl}) file. The program is invoked from the command line:

\begin{verbatim}
python archimate_to_eakg.py <input.xml> [output.ttl]
\end{verbatim}

If no output path is given, the result is written to \texttt{eakg\_output.ttl} in the same directory as the input file.

The translation pipeline is organized into three sequential passes over the XML document, each implementing one of the core translation rules defined in \autoref{chap:translation}.

% ------------------------------------------------------------------
\subsection{Namespaces and Lookup Tables}
\label{subsec:namespaces}

The program begins by declaring the RDF namespaces used throughout the graph. Three namespaces are central to the output:

\begin{itemize}
    \item \texttt{ex:} (\texttt{http://www.example.org/eakg\#}) --- used for all individual elements, relationships, and property predicates that originate from the concrete ArchiMate model.
    \item \texttt{archimate:} (\texttt{http://purl.org/eakg/archimate\#}) --- the ArchiMate metamodel vocabulary represented by the ontology described in \autoref{sec:ontology}, published at a persistent URL. Classes such as \texttt{archimate:BusinessActor}, relationship types such as \texttt{archimate:Association}, and schema-fixed attributes such as \texttt{archimate:accessType} are defined under this namespace.
    \item \texttt{smells:} (\texttt{http://purl.org/eakg/smells\#}) --- the EA smell ontology, also published at a persistent URL. Defines the \texttt{smells:hasSmell} predicate and seven smell individuals used during the enrichment phase.
\end{itemize}

A lookup table maps ArchiMate primitive type strings to their corresponding XSD datatype URIs. This table is used during Pass~1 to assign the correct \texttt{rdfs:range} to each property definition:

\begin{verbatim}
  "string"  -> xsd:string
  "boolean" -> xsd:boolean
  "integer" -> xsd:integer
  "float"   -> xsd:float
  "date"    -> xsd:date
  "time"    -> xsd:time
  "number"  -> xsd:decimal
\end{verbatim}

% ------------------------------------------------------------------
\subsection{Utility Functions}
\label{subsec:utilities}

Several helper functions are shared across passes:

\begin{itemize}
    \item \texttt{id\_to\_iri} and \texttt{property\_name\_to\_iri} create stable, sanitized IRIs; \texttt{resolve\_archimate\_class} maps XML types to ontology classes; and \texttt{emit\_lang\_literals} preserves language-tagged names and documentation. Detailed behavior is summarized in Annex~B.
\end{itemize}

% ------------------------------------------------------------------
\subsection{Pass 1: Property Definitions (Rule 3 --- Definition Half)}
\label{subsec:pass1}

The first pass registers property predicates and their XSD ranges, returning the lookup dictionary needed by the element and relationship passes.

% ------------------------------------------------------------------
\subsection{Pass 2: Elements (Rule 1)}
\label{subsec:pass2}

The second pass converts each XML element into a typed named individual, emits multilingual labels and comments, and writes its property values as direct triples using the Pass~1 dictionary.

% ------------------------------------------------------------------
\subsection{Pass 3: Relationships (Rule 2)}
\label{subsec:pass3}

The third pass turns each relationship identifier into an object-property predicate, links it to its ArchiMate relationship class, and emits the source-to-target triple plus any metadata and property values.

% ------------------------------------------------------------------
\subsection{Knowledge Graph Deployment and Usage}
\label{subsec:kg-deployment-usage}

The resulting KG then loaded in the triple store (importing the turtle file in GraphDB), and also import the Archimate Ontology which includes definition and hierarchy of used classes/predicate. Then the graph can be visualized or queried using SPARQL.

% ------------------------------------------------------------------
\section{Enrichment: Ontology Loading and Smell Tagging}
\label{sec:enrichment}

Enrichment is performed directly in the triplestore after loading the created graph.

\subsection{Ontology Loading}
\label{subsec:ontology-loading}

The enrichment ontology EA smell ontology in our case is loaded into the triplestore via its persistent URL \texttt{http://purl.org/eakg/smells}), alongside the translated model. This allow using resources and predicate defined in the enrichment ontology.

\subsection{EA Smell Tagging}
\label{subsec:smell-tagging}

Once the ontologies and model are loaded and RDFS reasoning is active, EA smells are detected and tagged using SPARQL queries. For each smell, the process follows two steps:

\begin{enumerate}
    \item A \textbf{detection \texttt{SELECT}} query identifies the affected elements or relations.
    \item A \textbf{tagging \texttt{INSERT}} query materializes a \texttt{smells:hasSmell} triple on each affected resource, using the same detection logic as its \texttt{WHERE} clause.
\end{enumerate}

The detection queries are unchanged from the analytical queries presented in Chapter~\ref{chap:queries}. The tagging queries are presented alongside them. Once tagged, affected resources can be retrieved with a single generic query by specifying the smell URI, and tags can be selectively reset.

% ------------------------------------------------------------------
\section{Tools and Technology Stack}
\label{sec:tools}

The implementation relies on the following tools and technologies:

\begin{itemize}
    \item \textbf{Python 3} as the programming language for both the translator and the enrichment module.
    \item \textbf{rdflib} (Python library) for constructing, manipulating, and serializing RDF graphs \parencite{rdflib2024}. It provides native support for Turtle, N-Triples, and other RDF serialization formats, as well as namespace management and literal typing.
    \item \textbf{xml.etree.ElementTree} (Python standard library) for parsing the ArchiMate Model Exchange File Format XML documents.
    \item \textbf{GraphDB} (Ontotext) as the target triplestore for deploying the resulting EAKG \parencite{ontotext2024graphdb}. GraphDB supports SPARQL 1.1 querying and RDFS reasoning, which enables the sub-property inference chains relied upon by our translation rules.
    \item \textbf{SPARQL 1.1} as the query language for analyzing the deployed knowledge graph.
    \item \textbf{RDFS reasoning} for inferring implicit relationships through the \texttt{rdfs:subPropertyOf} and \texttt{rdfs:subClassOf} hierarchies defined in the base ontology. This allows queries to match, for example, all \texttt{archimate:DependencyRelationship} instances without enumerating every concrete sub-property.
\end{itemize}
